% Generated draft: requires mathematical and source review.
\documentclass[11pt]{article}
\usepackage[margin=25mm]{geometry}
\usepackage{fontspec}
\setmainfont{Latin Modern Roman}
\usepackage{amsmath,amssymb,mathtools,unicode-math}
\setmathfont{Latin Modern Math}
\usepackage{enumitem,needspace,xurl,hyperref}
\setlength{\parindent}{0pt}
\setlength{\parskip}{5pt}
\setlength{\emergencystretch}{3em}
\newcommand{\N}{\mathbb{N}}
\newcommand{\Z}{\mathbb{Z}}
\newcommand{\Q}{\mathbb{Q}}
\newcommand{\R}{\mathbb{R}}
\newcommand{\C}{\mathbb{C}}
\newcommand{\F}{\mathbb{F}}
\newcommand{\A}{\mathbb{A}}
\newcommand{\PP}{\mathbb{P}}
\newcommand{\E}{\mathbb{E}}
\newcommand{\eps}{\varepsilon}
\newcommand{\dd}{\,\mathrm{d}}
\begin{document}
\section*{UnsolvedMath ID: 30003010}
\textbf{Source problem code:} catalogue entry
\subsection{Definitions and mathematical statement}
Let \(n,k\in\Z\) satisfy \(n\geq3\) and \(2\leq k\leq n-1\). Use the customary convention that a convex body in \(\R^n\) is a compact convex set with nonempty interior. Let \(G(n,k)\) denote the set of all \(k\)-dimensional linear subspaces of \(\R^n\). For \(H\in G(n,k)\), let \(P_H:\R^n\to H\) be orthogonal projection and define
\[
SO(H)=\{R:H\to H: R\text{ is linear and orthogonal},\ \det R=1\}.
\]
This interprets the source notation \(SO(k,H)\).

Interpreting directly congruent projections as congruence by orientation-preserving rigid motions within each projection subspace, the question is whether the following implication holds for every pair of convex bodies \(K,L\subseteq\R^n\):
\[
\left[\forall H\in G(n,k)\ \exists R_H\in SO(H)\ \exists t_H\in H:\quad P_HK=R_H(P_HL)+t_H\right]
\Longrightarrow
\left[\exists a\in\R^n\ \exists\varepsilon\in\{-1,1\}:\quad K=a+\varepsilon L\right].
\]
The rotations and projection translations may depend on \(H\); the conclusion requires one translation and one sign for the entire bodies. Here \(-L=\{-x:x\in L\}\).

There is an unresolved qualification in the supplied wording: directly congruent ordinarily permits translations, and the background refers to a direct-rigid-motion problem, but the statement also says that the congruences are induced by elements of \(SO(k,H)\), which consists of linear rotations. If that phrase excludes translations, the alternative hypothesis is \(\forall H\in G(n,k)\ \exists R_H\in SO(H):P_HK=R_H(P_HL)\). The supplied record does not explicitly settle this distinction. The convex-body convention above is also not explicitly defined in the source record.
\subsection{Short English statement}
For each \(2\leq k\leq n-1\), suppose every pair of corresponding orthogonal \(k\)-dimensional projections of two convex bodies is directly congruent. Must the bodies differ by a translation, possibly followed by central reflection? Whether translations are allowed in the projection congruences requires source review.
\subsection{Sources}
\begin{itemize}
\item Supplied local UnsolvedMath registry record, permanent ID \(30003010\), title: Reconstruction from Directly Congruent Convex Projections; collection: Oberwolfach Reports Open Problems; exact source problem code: \texttt{catalogue entry}. The record lists proposed year \(2015\) and identifies no proposer.
\item Supplied report citation: Convex Geometry and its Applications (2016), Oberwolfach Reports, pp.\ 3179--3269, DOI \texttt{10.4171/owr/2015/56}. \url{https://doi.org/10.4171/owr/2015/56}. The background describes these reports as documenting mathematical workshops and their open-problem discussions. The supplied citation is authoritative report provenance, but its contents have not been independently consulted here.
\item Supplied literature reference: Alfonseca, Cordier, and Ryabogin. \url{https://arxiv.org/abs/1412.2727}. No paper title or theorem-level citation is supplied.
\item Literature assessment checked at \texttt{2026-08-21}: supplied dated curation, not independently verified. The record labels the problem partially solved and reports affirmative results for the special-orthogonal version under restrictions such as two-dimensional projections without rotational symmetries, while describing the full direct-rigid-motion problem as open. The background explicitly calls this a dated literature triage derived from supplied OWR curation and requests independent verification of time-sensitive status and cited claims.
\item The registry labels the statement \texttt{corrected\_verified}; its accompanying verification note states that only safe notation, typography, and TeX-delimiter cleanup was applied, preserving hypotheses, quantifiers, and subquestions. This is supplied metadata, not independent mathematical verification of this draft.
\end{itemize}
\end{document}
